\subsection{Reminders concerning continuous linear representations}

Let G be a topological group. One recalls that a continuous linear representation of G in a topological K-vector space E is a linear representation \(\varrho\) of G in E such that the map from G × E to E defined by \((g, x) \mapsto \varrho(g)x\) is continuous (INT, VIII, p. 128, § 2, n° 1, def. 1). This map defines an action of G on E and the image of \(\varrho\) is contained in \(\mathcal{L}(\mathrm{E})\) . One says that a continuous representation \(\varrho\) is bounded if its image is bounded in the space \(\mathcal{L}(\mathrm{E})\) endowed with the topology of bounded convergence.

Remark. --- If K = R and if \(\varrho\) is a continuous linear representation of G in an R-vector space E, then the map \(\varrho_{(\mathbf{C})}: g \mapsto 1 \otimes \varrho(g)\) is a continuous linear representation of G in \(\mathrm{E}_{(\mathbf{C})}\) ; if \(\varrho\) is bounded, then \(\varrho_{(\mathbf{C})}\) is bounded.

For every topological K-vector space E, the homomorphism which to every \(g \in G\) associates \(1_{E}\) is a continuous linear representation of G in E, called the trivial representation of G in E.

Let H be a subgroup of G. The restriction of \(\varrho\) to H is a continuous linear representation of H in E, denoted \(\mathrm{Res}_{\mathrm{H}}^{\mathrm{G}}(\varrho)\) .

Let \(\varrho_{1}\) and \(\varrho_{2}\) of continuous linear representations of G in K-vector topological spaces E \(_{1}\) and E \(_{2}\) . A G-morphism \(u\) of \(\varrho_{1}\) in \(\varrho_{2}\) is a morphism of linear representations that is continuous, that is, a continuous linear map \(u\colon\) E \(_{1} \to\) E \(_{2}\) such that \(u \circ \varrho_{1}(g) = \varrho_{2}(g) \circ u\) for all \(g \in\) G. We denote by Hom \(_{G}\) ( \(\varrho_{1}, \varrho_{2}\) ) the vector space of G-morphisms from \(\varrho_{1}\) in \(\varrho_{2}\) . It is a closed subspace of the space \(\mathcal{L}\) (E \(_{1}\) ; E \(_{2}\) ) equipped with the topology of pointwise convergence.

Let \(\varrho\) a continuous linear representation of a group G in a K-vector topological space E. The identity map \(1_{\mathrm{E}}\) of E is a G-morphism of \(\varrho\) in \(\varrho\) , which is denoted \(1_{\varrho}\) . If \(\varrho_1\) , \(\varrho_2\) and \(\varrho_3\) are continuous linear representations of G into K-topological vector spaces \(\mathrm{E}_1\) , \(\mathrm{E}_2\) , \(\mathrm{E}_3\) respectively, and if \(u: \mathrm{E}_1 \to \mathrm{E}_2\) and \(v: \mathrm{E}_2 \to \mathrm{E}_3\) are G-morphisms, then \(v \circ u\) is a G-morphism.

A G-morphism \(u\) of \(\varrho_{1}\) in \(\varrho_{2}\) is a G-isomorphism if there exists a G-morphism \(v\) of \(\varrho_{2}\) in \(\varrho_{1}\) such that \(v\circ u=1_{\varrho_{1}}\) and \(u\circ v=1_{\varrho_{2}}\) It is necessary and sufficient for this that \(u\) be a G-morphism and an isomorphism of topological vector spaces of the space \(\varrho_{1}\) in the space of \(\varrho_{2}\) ; its inverse \(v=u^{-1}\) is indeed then a G-morphism. If there exists a G-isomorphism from \(\varrho_{1}\) in \(\varrho_{2}\) we say that these representations are isomorphic.

DEFINITION 1. --- Let \(\varrho\) a continuous linear representation of a topological group G in a topological vector space E. \(\pi\) a subrepresentation of \(\varrho\) If F is a closed subspace of E and if, for every \(g \in G\) the space F is stable under \(\varrho(g)\) and \(\pi(g)\) is the endomorphism of F induced by \(\varrho(g)\) by passage to the subspaces.

Let \(\varrho\) a continuous linear representation of a topological group G in a topological vector space E. The closure of \(\{0\}\) in E is a subrepresentation of \(\varrho\) More generally, the closure of a subspace stable under the endomorphisms \(\varrho(g)\) is a subrepresentation of \(\varrho\) .

A subrepresentation \(\pi\) of \(\varrho\) is uniquely determined by a closed subspace F, stable under all endomorphisms \(\varrho(g)\) This last condition will often be stated in the form “F is a G-invariant subspace of E.” One will then also say that F defines a subrepresentation of \(\varrho\) or sometimes, by abuse of language, that F is a subrepresentation of \(\varrho\) , or of E.

Let F be a subspace of E defining a subrepresentation \(\pi\) of \(\varrho\) . For every \(g \in G\) , the linear map \(\varrho(g)\) defines, by passing to the quotient, an endomorphism \(\widetilde{\varrho}(g)\) of E/F. The map \(g \mapsto \widetilde{\varrho}(g)\) is a continuous linear representation of G in E/F, and the canonical projection of E onto E/F is a G-morphism. We say that \(\widetilde{\varrho}\) is the quotient representation of G on E/F; it is also denoted \(\varrho/\pi\) .

The subspace of elements of E invariant under the action of G is a trivial subrepresentation of \(\varrho\) . We denote it by \(\varrho^{G}\) or \(E^{G}\) .

Let A be a subset of E. The closed subspace F generated by the elements \(\varrho(g)x\) , where \(g \in G\) and \(x \in A\) is a subrepresentation of \(\varrho\) called the subrepresentation of \(\varrho\) generated by A; if F = E, one says that A generates \(\varrho\) . If the subrepresentation F is finite-dimensional, one says that the subset A is G-finite.

Suppose that A is reduced to a single element \(x \in E\) . If A generates \(\varrho\) , one says that x is a cyclic vector of \(\varrho\) and that \(\varrho\) is a cyclic representation; one says that x is a G-finite vector if A is G-finite.

The set of G-finite vectors of \(\varrho\) is a vector subspace of E stable under \(\varrho\) ; it is not necessarily closed.

Let \((\varrho_{i})_{i\in\mathrm{I}}\) a family of continuous linear representations of G in locally convex topological K-vector spaces \((\mathrm{E}_{i})_{i\in\mathrm{I}}\) . Let E be the direct sum space of the spaces \(E_{i}\) . The map \(g\mapsto(\varrho_{i}(g))\) is a linear representation of G in E, called the sum, or direct sum, of the representations \(\varrho_{i}\) If I is finite, it is continuous.

If the representations \(\varrho_{i}\) are equal to a representation \(\varrho\) for all \(i \in I\) , we say that the direct sum of the representations \(\varrho_{i}\) is the sum of Card(I) copies of the representation \(\varrho\) , and it is also denoted \(\varrho^{\mathrm{Card(I)}}\) , or Card(I) \(\varrho\) .

Let \((\varrho_{i})_{i\in\mathrm{I}}\) and \((\pi_{j})_{j\in\mathrm{J}}\) of finite families of continuous linear representations of G in topological K-vector spaces \((\mathrm{E}_{i})_{i\in\mathrm{I}}\) and \((\mathrm{F}_{j})_{j\in\mathrm{J}}\) , respectively. The canonical isomorphism of K-vector spaces

\[
\operatorname{Hom} _ {\mathrm{K}} \left(\bigoplus_ {i \in \mathrm{I}} \mathrm{E} _ {i}, \bigoplus_ {j \in \mathrm{J}} \mathrm{F} _ {j}\right)\rightarrow \bigoplus_ {(i, j) \in \mathrm{I} \times \mathrm{J}} \operatorname{Hom} _ {\mathrm{K}} (\mathrm{E} _ {i}, \mathrm{F} _ {j})
\]

(A, II, p. 13, cor. 1) induces, by passage to subspaces, an isomorphism

\[
\mathrm{Hom} _ {\mathrm{G}} \Bigl (\bigoplus_ {i \in \mathrm{I}} \varrho_ {i}, \bigoplus_ {j \in \mathrm{J}} \pi_ {j} \Bigr) \to \bigoplus_ {(i, j) \in \mathrm{I} \times \mathrm{J}} \mathrm{Hom} _ {\mathrm{G}} (\varrho_ {i}, \pi_ {j})\tag{1}
\]

which is likewise said to be canonical.

Let \(\varrho_{1}\) and \(\varrho_{2}\) of linear representations of G in K-vector spaces E \(_{1}\) and E \(_{2}\) . For \(u \in \mathcal{L}(\mathrm{E}_{1};\mathrm{E}_{2})\) and \(g \in \mathrm{G}\) we set \(\varrho(g)u = \varrho_{2}(g) \circ u \circ \varrho_{1}(g^{-1})\) . The map \(g \mapsto \varrho(g)\) is a linear representation of G in \(\mathcal{L}(\mathrm{E}_{1};\mathrm{E}_{2})\) . The space of invariant elements of this representation coincides with \(\mathrm{Hom}_{\mathrm{G}}(\varrho_{1},\varrho_{2})\) .

Let \(\varrho\) a continuous linear representation of G in a locally convex space E. Recall (INT, VIII, p. 131, § 2, no. 2) that the contragredient representation \(\breve{\varrho}\) of \(\varrho\) is the linear representation of G in the dual E' of E defined by \(\breve{\varrho}(g) = {}^{t} \varrho(g^{-1})\) .

